\documentclass[12pt]{extarticle}

% ============================================================
% PACKAGES
% ============================================================

% Encodage et langue
\usepackage[french]{babel}

% Mise en page
\usepackage[a4paper, margin=2cm]{geometry}

\usepackage[francais,bloc,completemulti]{automultiplechoice}    

% Mathématiques
\usepackage{amsmath}
\usepackage{amssymb}

% Dessins
\usepackage{tikz}

% Listes
\usepackage{enumitem}

% Boîtes
\usepackage[most]{tcolorbox}

% Images
\usepackage{graphicx}

% Mise en page avancée
\usepackage{fancyhdr}
\usepackage{multicol}

% Tableaux
\usepackage{booktabs}
\usepackage{tabularx}
\usepackage{colortbl}

\usepackage[normalem]{ulem}

\usepackage{pgfplots}
\pgfplotsset{compat=1.18}
\usepackage{pgf-pie}


\usepackage{ifthen}

% ============================================================
% COULEURS
% ============================================================

% ------------------------------------------------------------
% Couleurs principales
% ------------------------------------------------------------

\definecolor{Bleu}{RGB}{40, 100, 180}
\definecolor{Vert}{RGB}{50, 150, 90}
\definecolor{Orange}{RGB}{230, 140, 40}

\definecolor{BleuClair}{RGB}{235, 245, 255}
\definecolor{VertClair}{RGB}{235, 250, 240}
\definecolor{OrangeClair}{RGB}{255, 245, 225}

% ------------------------------------------------------------
% Couleurs complémentaires
% ------------------------------------------------------------

% Rappels de cours — violet
\definecolor{Violet}{RGB}{125, 85, 170}
\definecolor{VioletClair}{RGB}{245, 238, 252}

% Théorèmes — rouge
\definecolor{Rouge}{RGB}{190, 55, 55}
\definecolor{RougeClair}{RGB}{255, 238, 238}

% Lemme — indigo
\definecolor{Indigo}{RGB}{75, 85, 155}
\definecolor{IndigoClair}{RGB}{240, 241, 252}

% Points méthode — turquoise
\definecolor{Turquoise}{RGB}{35, 145, 150}
\definecolor{TurquoiseClair}{RGB}{235, 250, 250}

% Erreurs / attention — jaune
\definecolor{Jaune}{RGB}{210, 165, 35}
\definecolor{JauneClair}{RGB}{255, 250, 225}

% ------------------------------------------------------------
% Neutres
% ------------------------------------------------------------

\definecolor{GrisClair}{RGB}{245, 245, 245}
\definecolor{Gris}{RGB}{100, 100, 100}


% ============================================================
% COMMANDES PERSONNALISÉES
% ============================================================

\newcommand{\bb}[1]{\textbf{#1}}
\newcommand{\ii}[1]{\textit{#1}}

\newcommand{\cm}{\,\mathrm{cm}}
\newcommand{\m}{\,\mathrm{m}}

\newcommand{\reponse}[1]{\boxed{#1}}

% ============================================================
% TABLEAUX
% ============================================================

\newcommand{\ligneentete}{%
    \rowcolor{Bleu}%
    \color{white}\bfseries
}

\newcommand{\entete}[1]{%
    \cellcolor{Bleu}\color{white}\bb{#1}%
}


\NewDocumentCommand{\hl}{O{JauneClair}m}{%
  \tcbox[
    on line,
    colback=#1,
    colframe=#1,
    boxrule=0pt,
    arc=4pt,
    left=0pt,
    right=0pt,
    top=1pt,
    bottom=1pt
  ]{#2}%
}

\NewDocumentCommand{\hlb}{O{JauneClair}m}{%
  \tcbox[
    on line,
    colback=#1,
    colframe=#1,
    boxrule=0pt,
    arc=4pt,
    left=0pt,
    right=0pt,
    top=1pt,
    bottom=1pt
    ]{\bb{#2}}%
}
\newlength{\answerwidth}


\newcommand{\answer}[2][VertClair]{%
    \ifmmode
        \settowidth{\answerwidth}{$\displaystyle #2$}%
        \ifanswers
            {\setlength{\fboxsep}{0pt}%
             \colorbox{#1}{$\displaystyle #2$}}%
        \else
            \makebox[0pt][l]{%
                \makebox[\answerwidth][l]{%
                    \leaders\hbox to 0.5em{\hss.\hss}\hfill
                }%
            }%
            \phantom{#2}%
        \fi
    \else
        \settowidth{\answerwidth}{#2}%
        \ifanswers
            \hl[#1]{#2}%
        \else
            \makebox[\answerwidth][l]{%
                \leaders\hbox to 0.5em{\hss.\hss}\hfill
            }%
        \fi
    \fi
}


\newcommand{\graphanswer}[1]{%
    \ifanswers
        #1%
    \else
        0%
    \fi
}

\newcommand{\exo}[1]{%
    \includegraphics[
        width=\linewidth,
        keepaspectratio
    ]{#1}%
}

\newcommand{\nodata}{}
\newcommand{\correctanswer}[1]{
    \explain{
            \newline\vspace{10pt}
            #1
    }
}

\newlength{\mylen}
\newtcolorbox{exercicebox}[1][\unskip]{
    colback=white,
    colframe=Bleu,
    boxrule=1pt,
    arc=3mm,
    title=\bb{Exercice #1}\AMCnumero{1},
    left=2mm,
    right=2mm,
    top=2mm,
    bottom=0mm
}

\newcounter{theorem}
\newtcolorbox{theoremebox}[1][\unskip]{
    colback=RougeClair,
    colframe=Rouge,
    boxrule=1pt,
    arc=3mm,
    title=\bb{Théorème \stepcounter{theorem}\thetheorem}\hspace{5pt}#1,
    left=5mm,
    right=5mm,
    top=4mm,
    bottom=4mm
}

\newcounter{probleme}
\newtcolorbox{probbox}[1][\unskip]{
    colback=RougeClair,
    colframe=Rouge,
    boxrule=1pt,
    arc=3mm,
    title=\bb{Problème \stepcounter{probleme}\theprobleme #1}\AMCnumero{1},
    left=5mm,
    right=5mm,
    top=2mm,
    bottom=2mm
}

\newenvironment{tableau}[1]
{%
    \renewcommand{\arraystretch}{1.3}%
    \begin{center}%
    \begin{tabularx}{\linewidth}{#1}%
    \toprule
}
{%
    \bottomrule
    \end{tabularx}%
    \end{center}%
    \renewcommand{\arraystretch}{1}%
}



% Set page layout after the config files
\geometry{
  paper=a4paper,
  hmargin=1cm,
  headheight=1cm,
  headsep=.3cm,
  footskip=0cm,
  top=1.5cm,
  left=1.3cm,
  right=1.3cm,
  bottom=1.5cm
}

%\geometry{hmargin=1cm,headheight=1cm,headsep=.3cm,footskip=0cm,top=1.5cm,bottom=1.5cm}

\def\AMCotextReserved{\emph{(Reservé)}}
\AMCopenOpts{boxframerule=0pt,boxframerulecol=red,framerule=0pt}
\newdimen\AMCpostOquest\AMCpostOquest=-7mm
\AMCformHSpace=.3em
\def\AMCbeginQuestion#1#2{\textbf{(#1)\hspace{2mm}}}

\begin{document}
\exemplaire{30}{    

% ---- Header ------
\fbox{
    \begin{minipage}{.43\linewidth}
        \namefield{Nom: }
        Prénom :\\
        Date:
    \end{minipage}
}
\hfill
\begin{minipage}{.48\linewidth}
    \AMCcodeGridInt[vertical=false]{etu}{2}
\end{minipage}
% ---- Fin Header ------

\begin{exercicebox}[1 \hspace{5mm}\it{Statistiques}]
        \begin{question}{qcmfrequence}\bareme{b=1}  
            Sur 40 élèves, 10 font du sport en club. La fréquence vaut…
            \begin{reponseshoriz}
                \mauvaise{0,10}
                \bonne{0,25}
                \mauvaise{4,0}
                \mauvaise{10}
                \mauvaise{0,40}
            \end{reponseshoriz}
        \end{question}


        \begin{question}{qcmmoyenne}\bareme{b=1}  
            La moyenne de 7 ; 9 ; 14 ; 10 est…
            \begin{reponseshoriz}
                \bonne{10}
                \mauvaise{9}
                \mauvaise{40}
                \mauvaise{11}
                \mauvaise{23}
            \end{reponseshoriz}
        \end{question}

        \begin{question}{qcmmoyennesomme}\bareme{b=1}  
            Une série de 6 notes a pour moyenne 12. La somme des notes vaut…
            \begin{reponseshoriz}
                \bonne{72}
                \mauvaise{12}
                \mauvaise{60}
                \mauvaise{18}
                \mauvaise{2}
            \end{reponseshoriz}
        \end{question}

        \begin{minipage}{0.4\textwidth}
            \includegraphics[width=\textwidth]{images/diagramme_barre.png}
        \end{minipage}
        \hfill
        \begin{minipage}{0.6\textwidth}

            \begin{question}{qcmdiagramme}\bareme{b=1}  
                Nombre de matchs joués dans la saison :
                \begin{reponseshoriz}
                    \bonne{28}
                    \mauvaise{9}
                    \mauvaise{6}
                    \mauvaise{7}
                \end{reponseshoriz}
            \end{question}

            \begin{question}{qcmdiagramme2}\bareme{b=1}  
                Nombre de matchs où l’équipe a marqué au moins 3 buts : 
                \begin{reponseshoriz}
                    \bonne{8}
                    \mauvaise{5}
                    \mauvaise{20}
                    \mauvaise{9}
                \end{reponseshoriz}
            \end{question}

        \end{minipage}

\end{exercicebox}

\begin{exercicebox}[2 \hspace{5mm}\it{Opérations}]

    \begin{question}{qcmpriodiv}\bareme{b=1}  
        \begin{minipage}{0.20\textwidth}
        $45\div9\div5$ vaut...
        \end{minipage}
        \hfill
        \begin{minipage}{0.7\textwidth}
        \begin{reponseshoriz}
            \mauvaise{25}
            \mauvaise{5}
            \bonne{1}
            \mauvaise{0,2}
        \end{reponseshoriz}
        \end{minipage}
    \end{question}

    \begin{question}{qcmpriomultiadd}\bareme{b=1}  
        \begin{minipage}{0.20\textwidth}
            $3\times4+2\times5$ vaut…
        \end{minipage}
        \hfill
        \begin{minipage}{0.7\textwidth}
            \begin{reponseshoriz}
                \mauvaise{70}
                \mauvaise{26}
                \mauvaise{50}
                \bonne{22}
            \end{reponseshoriz}
        \end{minipage}
    \end{question}

    \begin{question}{prez}    
          Dans un calcul sans parenthèses, on effectue d'abord…
          \begin{reponses}
              \bonne{les multiplications et divisions}
              \mauvaise{les additions et soustractions}
              \mauvaise{les opérations de gauche à droite}
              \mauvaise{l'opération avec le plus grand nombre}
          \end{reponses}
    \end{question}

    \begin{question}{qcmprioparentheses}\bareme{b=1}
        \begin{minipage}{0.20\textwidth}
            $(4+6)\times2$ vaut…
        \end{minipage}
        \hfill
        \begin{minipage}{0.60\textwidth}
            \begin{reponseshoriz}
                \mauvaise{16}
                \bonne{20}
                \mauvaise{14}
                \mauvaise{12}
            \end{reponseshoriz}
        \end{minipage}
    \end{question}


    \begin{question}{qcmvocab}\bareme{b=1}  
        La somme du quotient de 4 par 6 et de 3 s'écrit
        \begin{reponseshoriz}
            \bonne{$4\div6+3$}
            \mauvaise{$4\div(6+3)$}
            \mauvaise{$3+6\div4$}
            \mauvaise{$4\times(3+6)$}
        \end{reponseshoriz}
    \end{question}

\end{exercicebox}

}   
\end{document}
